\section{Introduction}

\label{sec:introduction}

Blockchain transparency is simultaneously its greatest feature and its most
significant limitation. Every transaction amount, every account balance, every
smart contract state variable is publicly visible. This transparency enables
trustless verification but prohibits use cases that require confidentiality:

\begin{itemize}[leftmargin=2em]
  \item \textbf{Dark pools}: Institutional trading requires hidden orderbooks. On
        transparent chains, every pending order is visible, enabling front-running.
  \item \textbf{Private voting}: On-chain governance with visible votes before
        tallying enables vote-buying and last-minute strategic voting.
  \item \textbf{Sealed-bid auctions}: Auction mechanisms (Vickrey, Dutch) require
        bid secrecy until the reveal phase. Transparent chains leak bids.
  \item \textbf{Regulatory compliance}: Regulation ATS (Alternative Trading Systems)
        requires that certain trading data remain confidential from non-participants.
\end{itemize}

Fully homomorphic encryption (FHE) resolves this tension by enabling computation on
encrypted data. A smart contract operating on FHE ciphertexts can add encrypted
balances, compare encrypted values, and update encrypted state---all without
decrypting. Only authorized parties (key holders) can decrypt the final results.

This paper describes the Lux FHE implementation: a native Go cryptographic backend
with NTT SIMD acceleration, integrated as an EVM precompile and deployed on the
LiquidFHE chain.

%% ========================================================================
%%  2. FHE PRIMER
%% ========================================================================
